%======================================================================
\section{A Bell-gap contrast and finite blocks}\label{sec:boundary}
%======================================================================

This section uses the Bell-gap contract of \cref{sec:prelim-usable} for its conditional $I_{3322}$ statement, then returns to
classes~(S) and~(C) for the finite-block observation. The two contracts are compared in \cref{sec:status}; whether Theorem~A
extends beyond classes~(S) and~(C) is open, and this section does not decide it.

\subsection{The \texorpdfstring{$I_{3322}$}{I3322} contrast}\label{sec:i3322}

We use~\cite{Dou26} in its arXiv version~2, whose mathematical content is that of version~4.1.0 of its archive. In its
Collins--Gisin normalization it proves finite-dimensional nonattainment of the quantum supremum $S$, certified to lie in
$(0.2508753845015185,0.250875388108398]$, for mixed states and POVMs~\cite[Theorem~1]{Dou26}; attainment by a normal spatial
strategy on $\ell^2(\Z)\otimes\ell^2(\Z)$~\cite[Theorem~2]{Dou26}; and $D(\varepsilon)=\Theta(\ln(1/\varepsilon))$ for the least local dimension at
Bell gap $\varepsilon=S-B$, with eventual upper coefficient $23.9010650$ in natural logarithms, not asserted optimal and with no
explicit onset~\cite[Theorem~6]{Dou26}. Nonattainment was proved independently by Pauwels~\cite[Theorem~2]{Pau26}, and nonattainment
together with attainment by a spatial strategy on separable infinite-dimensional Hilbert spaces by
Coladangelo~\cite[Theorem~11.1]{Col26}; the infinite-dimensional construction of the optimum, and the conjecture that finite
dimension does not reach it, are due to P\'al and V\'ertesi~\cite{PV10}. The tail estimate
we use is that of~\cite[Section~5.2]{Dou26}: the Schmidt coefficients of the spatial attainer constructed there decay exponentially
at both ends. Machine checks cover only parts of these works. Pauwels' Lean formalization covers pure states and projective
measurements, and the reduction from mixed states and POVMs lies outside it~\cite[Appendix~D]{Pau26}; the Lean checks
of~\cite{Dou26} cover selected numerical, scalar and finite accounting statements, and its analytic proofs, including spatial
attainment and the tail estimate, are not formalized. We carry that status.

\begin{theorem}[A conditional Bell-gap contrast]\label{thm:i3322}
Conditional on the spatial attainment theorem~\cite[Theorem~2]{Dou26} and the tail bound $\kappa_\pm\ge0.08367827985$
of~\cite[Section~5.2]{Dou26}, the infimum of the original one-way distillable entanglement, in either direction, over strategies with Bell gap at most $\varepsilon$ is $O(1)$ as $\varepsilon\to0$, although no finite-dimensional strategy
attains $S$ and the least dimension grows like $\ln(1/\varepsilon)$.
\end{theorem}

\begin{proof}
The normal pure attainer of~\cite{Dou26} is $\sum_j\lambda_j|j,j\rangle$ with geometric tails: for each small $\eta>0$,
$\lambda_j^2\le A_\eta\exp[-(\kappa_\pm-\eta)|j|]$ on the two tails, with $\kappa_\pm\ge0.08367827985$ per index~\cite[Section~5.2]{Dou26}.
So $p_j=\lambda_j^2$ has finite Shannon entropy $H(p)$; finite norm alone would not give this. The near-optimal finite strategies
of~\cite[Section~5.2]{Dou26} are the renormalized windows of this same Schmidt vector on finite intervals $I$, with
$Z_I=\sum_{j\in I}p_j\to1$, and their Bell gaps tend to zero. Their Schmidt distributions $q_j=p_j/Z_I$ satisfy
\[
H(q)=Z_I^{-1}\sum_{j\in I}(-p_j\log p_j)+\log Z_I\le H(p)/Z_I,
\]
since every term is nonnegative and $\log Z_I\le0$. These states are pure, so both directional rates equal $H(q)$, which stays bounded.
\end{proof}

No matching positive lower bound, exact limit or minimizing state is asserted, and the logarithmic dimension theorem does not
force any growth of $\ED$. The contract cannot be changed by rescaling: the nonsignaling value of $I_{3322}$ is one, above the
quantum optimum $S\simeq0.2509$. With the joint coefficients $\bigl(\begin{smallmatrix}1&1&1\\1&1&-1\\1&-1&0\end{smallmatrix}\bigr)$ and
marginal coefficients $(-1,0,0)$ and $(-2,-1,0)$, the behaviour with all marginals $\frac12$ and $P(1,1)=\frac12$ on the positive
entries and $0$ on the negative ones has value $-2+\frac62=1$, so an affine rescaling $aB+b$ with quantum maximum one would score
$1+a(1-S)>1$ on it, impossible for a probability of winning.

\subsection{Normal perfect strategies have finite perfect blocks}\label{sec:normal-perfect}

\begin{remark}[Finite blocks; known]\label{prop:normal-perfect}
In either class of \cref{def:law-classes}, any normal spatial perfect strategy, even with infinite-dimensional ambient systems and an
arbitrary normal mixed input, yields a finite-dimensional perfect strategy.
\end{remark}

\begin{proof}
For the synchronous class this is the theorem of Kim, Paulsen and Schafhauser that $C^s_q(n,m)=C^s_{qs}(n,m)$~\cite[Theorem~3.10 and
Corollary~3.11]{KPS18}. Directly: after local projectivization, perfect diagonal consistency gives $P^A_{xa}\Psi=Q^B_{xa}\Psi$ on a
purification; taking partial traces, the trace-class Bob marginal commutes with every Bob projection. A nonzero eigenvalue has a
finite-dimensional eigenspace invariant under these projections, and the losing products, which annihilate the support (move one
factor to Alice), vanish on it; the transpose maximally entangled realization of this block is a finite perfect strategy. For a
native constraint-variable game the same argument applies with the incidence mirrors and Alice's commuting legal row
projections, which is the standard operator-solution reconstruction~\cite{CM14}.
\end{proof}

So a game in the classes of Theorem~A without a finite perfect strategy has no normal perfect strategy at all, of any entropy.
Normal infinite-dimensional correlations that have no finite-dimensional realization were constructed by Coladangelo and
Stark~\cite{CoSt20}; the finite-block conclusion has a more restricted scope, perfect strategies of games in the two classes.
The $I_{3322}$ attainer is normal and of finite entropy, so a normality-preserving adapter that turned it into a perfect strategy
of a fixed target game in the classes of Theorem~A without a finite-dimensional perfect strategy would give that target a normal
spatial perfect strategy, which it does not have. For such a target, any sequence of finite-dimensional near-perfect strategies has
diverging total one-way distillable entanglement, by Theorem~A. A claim about the additional resources an adapter needs would
require a specified conversion contract, which we do not give.
